\section*{16. 逆序对的分治算法}

给定实数序列A，若i<j且A[i]>A[j]，则(i,j)为逆序对。要求用分治方法求序列中逆序对的个数。

\begin{itemize}
    \item \textbf{逆序对定义}：前大后小的元素对，如[3,1,2]的逆序对为(3,1)、(3,2)
    \item \textbf{分治分解}：将序列均分为左子序列L和右子序列R
    \item \textbf{递归求解}：递归计算L和R内部的逆序对个数countL和countR
    \item \textbf{合并统计}：归并L和R为有序序列，同时统计跨L、R的逆序对个数cross
    \item \textbf{跨区间统计}：归并时若L[i]>R[j]，则L[i]及L[i]之后的所有元素都与R[j]构成逆序对，count+=mid-i+1
    \item \textbf{时间复杂度}：归并排序的合并步骤嵌入统计，$T(n)=2T(n/2)+O(n)$，总计$O(n\log n)$
\end{itemize}

\begin{lstlisting}[language=C++, style=cppstyle]
int mergeAndCount(vector<int>& A, int low, int mid, int high) {
    vector<int> L(A.begin() + low, A.begin() + mid + 1);
    vector<int> R(A.begin() + mid + 1, A.begin() + high + 1);
    int i = 0, j = 0, k = low;
    int count = 0;
    while (i < L.size() && j < R.size()) {
        if (L[i] <= R[j]) {
            A[k++] = L[i++];
        } else {
            A[k++] = R[j++];
            count += L.size() - i;
        }    }}
    while (i < L.size()) A[k++] = L[i++];
    while (j < R.size()) A[k++] = R[j++];
    return count;}
int countInversions(vector<int>& A, int low, int high) {
    if (low >= high) return 0;
    int mid = (low + high) / 2;
    int countL = countInversions(A, low, mid);
    int countR = countInversions(A, mid + 1, high);
    int countCross = mergeAndCount(A, low, mid, high);
    return countL + countR + countCross;}}}
int countInversions(vector<int>& A) {
    return countInversions(A, 0, A.size() - 1);}}}
\end{lstlisting}

\begin{enumerate}
    \item 把序列从中间分成左右两半
    \item 递归统计左半内部的逆序对个数
    \item 递归统计右半内部的逆序对个数
    \item 归并左右两半为有序序列，同时统计跨区间的逆序对
    \item 统计跨区间时，如果左边元素大于右边，那左边剩下的元素也都大于右边当前元素
    \item 总逆序对数等于"左内+右内+跨区间"之和
\end{enumerate}
